% SEGMENT OLP-0536-S01
% Part: set-theory
% Chapter: story
% Section: grundgesetze

\documentclass[../../../include/open-logic-section]{subfiles}

\begin{document}

\olfileid{sth}{story}{blv}

\olsection{பின்னிணைப்பு: ஃப்ரேகேயின் அடிப்படை விதி V}

% SOURCE CORRECTION TA-STH-004: remove the source's possessive from “Russell's formulated his paradox.”
\olref[rus]{sec} இல், ஃப்ரேகே தமது \emph{Grundgesetze} நூலில்
விளக்கிய அமைப்பிற்கு ஒரு சிக்கலாக ரஸ்ஸல் தமது முரண்பாட்டை வகுத்தார்
என்று விளக்கினோம். ஃப்ரேகேயின் அமைப்பில் முறைசாராப் பண்புவழிக்
கணமாக்கலின் நேரடியான கூற்று இடம்பெறவில்லை. எனவே, ஃப்ரேகேயின் அமைப்பில்
எது \emph{இடம்பெற்றது}, அது முறைசாராப் பண்புவழிக் கணமாக்கலுடனும்
ரஸ்ஸலின் முரண்பாட்டுடனும் எவ்வாறு தொடர்புடையது என்பவற்றை இந்தப்
பின்னிணைப்பில் மிகச் சுருக்கமாக விளக்குவோம்.

ஃப்ரேகேயின் அமைப்பு \emph{இரண்டாம்வரிசையானது}; மேலும் அது ஒரு
\emph{கருத்துருவின் விரிவுநிலை} என்ற கருத்தாக்கத்தை வகுப்பதற்காக
வடிவமைக்கப்பட்டது.\footnote{கறாராகச் சொன்னால், ஃப்ரேகே இன்னும் பொதுவான
கருத்தாக்கம் ஒன்றை முறைசாராக்க முயல்கிறார்: ஒரு சார்பின்
“மதிப்பு-வீச்சு”. கருத்துருக்களின் விரிவுநிலைகள் அந்தப் பொதுவான
கருத்தாக்கத்தின் ஒரு சிறப்பு நிகழ்வு. விவரங்களுக்கு
\citet[pp.\ 8--9]{Heck2012} ஐப் பார்க்கவும்.} ஃப்ரேகேயால்
உந்தப்பட்ட குறியீட்டைப் பயன்படுத்தி, \emph{கருத்துரு~$F$ இன்
விரிவுநிலையை} $\fregeext{x}{F(x)}$ என எழுதுவோம். இந்த அமைப்பு,
“$F$” என்ற ஒரு \emph{பயனிலையை} எடுத்துக்கொண்டு, அதை
“$\fregeext{x}{F(x)}$” என்ற ஒரு (முதல்வரிசை) \emph{சொல்லாக}
மாற்றுகிறது. இந்த அமைப்பைப் பயன்படுத்தி, ஃப்ரேகே உறுப்பாதலுக்குப்
பின்வரும் \emph{வரையறையை} வழங்கினார்:
\[
	a \in b =_\text{df} \exists G(b = \fregeext{x}{G(x)} \land Ga)
\]
தோராயமாக: $a \in b$ ஆகும்போது, அப்போதும் அப்போது மட்டுமே, விரிவுநிலை
$b$ ஆக உள்ள ஒரு கருத்துருவின் கீழ் $a$ வருகிறது. (“$\exists G$” என்ற
அளவையடை இரண்டாம்வரிசையானது என்பதைக் கவனிக்கவும்.) \emph{அடிப்படை
விதி V} என்று அறியப்படும் பின்வரும் கொள்கையையும் ஃப்ரேகே ஏற்றார்:
$$\fregeext{x}{F(x)} = \fregeext{x}{G(x)} \liff \forall x (Fx \liff Gx)$$
தோராயமாக: கருத்துருக்கள் ஒரே பொருள்களுக்கு நிறைவானவையாக இருக்கும்போது,
அப்போதும் அப்போது மட்டுமே, அவை ஒரே விரிவுநிலைகளைக் கொண்டிருக்கும்.
(மீண்டும், “$F$” மற்றும் “$G$” இரண்டும் பயனிலை இடத்தில் உள்ளன.)
இப்போது ஓர் எளிய கொள்கை உறுப்பாதலை ஒரு பண்பு நிறைவுறுத்தப்படுவதுடன்
இணைக்கிறது:

\begin{lem}[\emph{Grundgesetze} நூலில்]\ollabel{lem:Fregeextensions}
$\forall F \forall a(a \in \fregeext{x}{F(x)} \liff Fa)$
\end{lem}

\begin{proof} 
$F$ மற்றும் $a$ ஐ நிர்ணயிக்க. இப்போது $a \in \fregeext{x}{F(x)}$
ஆகும்போது, அப்போதும் அப்போது மட்டுமே, $\exists G(\fregeext{x}{F(x)}
= \fregeext{x}{G(x)} \land Ga)$ (உறுப்பாதலின் வரையறையால்);
அப்போதும் அப்போது மட்டுமே, $\exists G(\forall x(Fx \liff Gx) \land Ga)$
(அடிப்படை விதி V ஆல்); அப்போதும் அப்போது மட்டுமே, $Fa$
(தொடக்கநிலை இரண்டாம்வரிசைத் தருக்கத்தால்).
\end{proof}

இதிலிருந்து முறைசாராப் பண்புவழிக் கணமாக்கல் கிட்டத்தட்ட உடனடியாகக்
கிடைக்கிறது:

\begin{lem}[\emph{Grundgesetze} நூலில்.]
$\forall F \exists s \forall a (a \in s \liff Fa)$
\end{lem}

\begin{proof}
$F$ ஐ நிர்ணயிக்க; இப்போது \olref{lem:Fregeextensions} ஆனது
$\forall a (a \in \fregeext{x}{F(x)} \liff Fa)$ ஐத் தருகிறது; எனவே
இருப்பளவைப் பொதுமைப்படுத்தலால் $\exists s\forall a(a \in s \liff Fa)$.
$F$ தன்னிச்சையாக இருந்ததால் முடிவு பின்வருகிறது.
\end{proof}

$\forall x(Fx \liff x \notin x)$ ஆல் $F$ தரப்படுவதாகக் கொள்வதன் மூலம்
ரஸ்ஸலின் முரண்பாடு பின்வருகிறது.

\end{document}
